% GENERATED FILE. Edit canonical lessons, then run npm run generate:books.
% canonical-source-hash: fnv1a64:37f8b831e51b3e0a
% canonical-lessons: BN-C158-koborsthan, BN-C158-prasad, BN-C158-durgo, BN-C158-minar, BN-C158-siling

\chapter{Graveyard, Palace, Fort, Minaret, Ceiling}
\label{ch:bn-graveyard-palace-fort-minaret-ceiling}
\hlchaptermodality{158}

\begin{chapteropening}
\textbf{By the end of this chapter:} \emph{I can say a graveyard, a palace, a fort, a minaret and a ceiling.}

Complete the last lesson of chapter 158: I can say a graveyard, a palace, a fort, a minaret and a ceiling.
\end{chapteropening}

\section[\texorpdfstring{\bn{কবরস্থান} (kôbôrsthān)}{কবরস্থান (kôbôrsthān)}]{\bn{কবরস্থান} (kôbôrsthān) — a graveyard}

\label{lesson:BN-C158-koborsthan}



\begin{quote}
Before the new one: say the Bengali for a canal, then the Bengali for a zoo.
\end{quote}

\subsection*{You'll want to know: \bn{কবরস্থান}}
\textbf{\bn{কবরস্থান}} — \emph{kôbôrsthān} — \textquotedblleft{}a graveyard\textquotedblright{}.

\begin{cousinweb}
\textbf{\bn{কবরস্থান}} (\emph{kôbôrsthān}) is \emph{kôbôr}, \textquotedblleft{}grave\textquotedblright{}, from Arabic \emph{qabr}, \textquotedblleft{}grave\textquotedblright{}, which reached Bengali through Persian (Lane), + \emph{sthān}, \textquotedblleft{}place\textquotedblright{}, Sanskrit \emph{sthāna} taken over whole. Persian has the matching \emph{qabristān}.
\end{cousinweb}

\begin{grammarlens}[title={What you've built}]
One more word for this chapter.
\end{grammarlens}

\subsection*{Your turn}
\begin{itemize}
\raggedright
  \item \emph{Say it:} \emph{kôbôrsthān}
  \item \emph{Say it:} \emph{kôbôrsthān}, once more
  \item \emph{Say it:} say \emph{kôbôrsthān}
  \item \emph{Recall:} say the Bengali for a canal, then the Bengali for a zoo, then say \emph{kôbôrsthān} again
\end{itemize}

\begin{tcolorbox}[breakable,colback=teal!4,colframe=teal!35!black,title={Before you move on}]
What does \bn{কবরস্থান} mean? (\textquotedblleft{}A graveyard\textquotedblright{}.) Say it once more.
\end{tcolorbox}
\section[\texorpdfstring{\bn{প্রাসাদ} (prāsād)}{প্রাসাদ (prāsād)}]{\bn{প্রাসাদ} (prāsād) — a palace}

\label{lesson:BN-C158-prasad}



\begin{quote}
Before the new one: say the Bengali for a zoo, then the Bengali for a graveyard.
\end{quote}

\subsection*{You'll want to know: \bn{প্রাসাদ}}
\textbf{\bn{প্রাসাদ}} — \emph{prāsād} — \textquotedblleft{}a palace\textquotedblright{}.

\begin{cousinweb}
\textbf{\bn{প্রাসাদ}} (\emph{prāsād}) is Sanskrit \emph{prāsāda}, \textquotedblleft{}lofty mansion\textquotedblright{}, taken over whole as a learned loan, a tatsama (CDIAL 8971).
\end{cousinweb}

\begin{grammarlens}[title={What you've built}]
One more word for this chapter.
\end{grammarlens}

\subsection*{Your turn}
\begin{itemize}
\raggedright
  \item \emph{Say it:} \emph{prāsād}
  \item \emph{Say it:} \emph{prāsād}, once more
  \item \emph{Say it:} say \emph{prāsād}
  \item \emph{Recall:} say the Bengali for a zoo, then the Bengali for a graveyard, then say \emph{prāsād} again
\end{itemize}

\begin{tcolorbox}[breakable,colback=teal!4,colframe=teal!35!black,title={Before you move on}]
What does \bn{প্রাসাদ} mean? (\textquotedblleft{}A palace\textquotedblright{}.) Say it once more.
\end{tcolorbox}
\section[\texorpdfstring{\bn{দুর্গ} (durgô)}{দুর্গ (durgô)}]{\bn{দুর্গ} (durgô) — a fort}

\label{lesson:BN-C158-durgo}



\begin{quote}
Before the new one: say the Bengali for a graveyard, then the Bengali for a palace.
\end{quote}

\subsection*{You'll want to know: \bn{দুর্গ}}
\textbf{\bn{দুর্গ}} — \emph{durgô} — \textquotedblleft{}a fort\textquotedblright{}.

\begin{cousinweb}
\textbf{\bn{দুর্গ}} (\emph{durgô}) is Sanskrit \emph{durga}, \textquotedblleft{}hard to reach; a fort\textquotedblright{}, taken over whole as a learned loan, a tatsama.
\end{cousinweb}

\begin{grammarlens}[title={What you've built}]
One more word for this chapter.
\end{grammarlens}

\subsection*{Your turn}
\begin{itemize}
\raggedright
  \item \emph{Say it:} \emph{durgô}
  \item \emph{Say it:} \emph{durgô}, once more
  \item \emph{Say it:} say \emph{durgô}
  \item \emph{Recall:} say the Bengali for a graveyard, then the Bengali for a palace, then say \emph{durgô} again
\end{itemize}

\begin{tcolorbox}[breakable,colback=teal!4,colframe=teal!35!black,title={Before you move on}]
What does \bn{দুর্গ} mean? (\textquotedblleft{}A fort\textquotedblright{}.) Say it once more.
\end{tcolorbox}
\section[\texorpdfstring{\bn{মিনার} (minār)}{মিনার (minār)}]{\bn{মিনার} (minār) — a minaret, a tower}

\label{lesson:BN-C158-minar}



\begin{quote}
Before the new one: say the Bengali for a palace, then the Bengali for a fort.
\end{quote}

\subsection*{You'll want to know: \bn{মিনার}}
\textbf{\bn{মিনার}} — \emph{minār} — \textquotedblleft{}a minaret, a tower\textquotedblright{}.

\begin{cousinweb}
\textbf{\bn{মিনার}} (\emph{minār}) is a loan from Arabic \emph{manār}, \textquotedblleft{}lighthouse, minaret\textquotedblright{}, which reached Bengali through Persian. Its Arabic root is \emph{n-w-r} (Lane). The word belongs with Arabic \emph{nūr}, \textquotedblleft{}light\textquotedblright{}.
\end{cousinweb}

\begin{grammarlens}[title={What you've built}]
One more word for this chapter.
\end{grammarlens}

\subsection*{Your turn}
\begin{itemize}
\raggedright
  \item \emph{Say it:} \emph{minār}
  \item \emph{Say it:} \emph{minār}, once more
  \item \emph{Say it:} say \emph{minār}
  \item \emph{Recall:} say the Bengali for a palace, then the Bengali for a fort, then say \emph{minār} again
\end{itemize}

\begin{tcolorbox}[breakable,colback=teal!4,colframe=teal!35!black,title={Before you move on}]
What does \bn{মিনার} mean? (\textquotedblleft{}A minaret, a tower\textquotedblright{}.) Say it once more.
\end{tcolorbox}
\section[\texorpdfstring{\bn{সিলিং} (siling)}{সিলিং (siling)}]{\bn{সিলিং} (siling) — a ceiling}

\label{lesson:BN-C158-siling}



\begin{quote}
Before the new one: say the Bengali for a fort, then the Bengali for a minaret.
\end{quote}

\subsection*{You'll want to know: \bn{সিলিং}}
\textbf{\bn{সিলিং}} — \emph{siling} — \textquotedblleft{}a ceiling\textquotedblright{}.

\begin{cousinweb}
\textbf{\bn{সিলিং}} (\emph{siling}) is a loan from English \emph{ceiling}.
\end{cousinweb}

\begin{grammarlens}[title={What you've built}]
One more word for this chapter.
\end{grammarlens}

\subsection*{Your turn}
\begin{itemize}
\raggedright
  \item \emph{Say it:} \emph{siling}
  \item \emph{Say it:} \emph{siling}, once more
  \item \emph{Say it:} say \emph{siling}
  \item \emph{Recall:} say the Bengali for a fort, then the Bengali for a minaret, then say \emph{siling} again
\end{itemize}

\begin{tcolorbox}[breakable,colback=teal!4,colframe=teal!35!black,title={Before you move on}]
What does \bn{সিলিং} mean? (\textquotedblleft{}A ceiling\textquotedblright{}.) Say it once more.
\end{tcolorbox}
